% =========================================================================
% experimental_setup.tex
% =========================================================================

\section{Experimental Setup}
\label{sec:setup}

\paragraph{Model.}
All inference experiments use GPT-2 Medium~\cite{radford2019gpt2}
(345M parameters, 24 layers, 16 attention heads, $\dhead = 64$,
maximum context length 1{,}024 tokens) loaded from the
\texttt{gpt2-medium} checkpoint on Hugging Face.

\paragraph{Hardware.}
A single NVIDIA RTX~4060 Laptop GPU (8\,GB VRAM) running CUDA~12.x.

\paragraph{Software.}
PyTorch~2.5.1, Hugging Face Transformers, SciPy (for Beta distribution
fitting in theorem validation).
All quantization is implemented in pure PyTorch without fused CUDA
or Triton kernels.

\paragraph{Datasets.}
\begin{itemize}[noitemsep,topsep=1pt,parsep=0pt,partopsep=0pt]
\item \textbf{WikiText-2}~\cite{merity2017pointer}: Used for strided
  perplexity evaluation with a 1{,}024-token sliding window.
  The published full-context perplexity for GPT-2 Medium is 19.69;
  the measured strided perplexity is 19.03.
  This gap reflects the difference between full-document evaluation and
  the 1{,}024-token strided window imposed by GPT-2's context limit.
  % Source: README.md line 235 (published = 19.69);
  %   results/final_comparison.json → Baseline FP16/wikitext_ppl = 19.03
\item \textbf{HellaSwag}~\cite{zellers2019hellaswag}: Zero-shot sentence
  completion accuracy (100~examples).
\item \textbf{NIAH (Needle-in-a-Haystack):} Natural-language
  question-answer retrieval at context lengths 256, 512, 768, and 1{,}024
  tokens with needle placement at Early, Middle, and Late positions.
\item \textbf{MQAR (Multi-Query Associative Recall):} Synthetic
  key-value association task at sequence lengths 64, 128, and 256 with
  4, 8, and 16 key-value pairs.
\item \textbf{LongBench-lite:} Passage QA, summarization
  (ROUGE-1), and code completion subtasks, truncated to 1{,}024 tokens.
\item \textbf{WikiText-103:} Used for training the 15M-parameter
  linear-attention MC-RNN models from scratch (20{,}000 steps).
\end{itemize}

\paragraph{Quantization schemes evaluated.}
This study compares six schemes across the benchmark suite:
\begin{enumerate}[noitemsep,topsep=1pt,parsep=0pt,partopsep=0pt]
\item \textbf{Baseline FP16:} No quantization (16-bit floating point).
\item \textbf{Original TurboQuant:} QR rotation + Lloyd-Max + QJL
  residual correction, at three-bit and four-bit.
\item \textbf{WHT + Asymmetric QJL (Adaptive):} WHT rotation +
  Lloyd-Max, QJL residual disabled at $\dhead = 64$, at three-bit and
  four-bit.
\item \textbf{MC-TurboQuant:} WHT + Asymmetric four-bit combined
  with GRM segmented caching.
\item \textbf{Outlier Channel Splitting:} Channel-aware bit allocation
  at 2.5-bit and 3.5-bit average rates.
\item \textbf{FibQuant:} Joint two-dimensional radial-angular codebook
  (tested but not included in the main comparison due to inferior
  performance; see Section~\ref{sec:method:fibquant}).
\end{enumerate}

\paragraph{Hyperparameters.}
All experiments use random seed~42.
Generation benchmarks use 50~tokens with greedy decoding.
The MC-RNN training uses a batch size of 16, learning rate $3 \times 10^{-4}$,
and the Adam optimizer for 20{,}000 steps on WikiText-103.
Segment sizes for Memory Caching experiments range from 8 to 256 tokens.
Nearest-neighbor evaluation uses 500 query vectors.

\paragraph{Reproducibility.}
All benchmark scripts, model weights (via Hugging Face), and result
JSON files are included in the repository.
Results can be reproduced by running the scripts in the order
documented in the repository README.
